% 分类目录：dl
\documentclass[12pt, a4paper]{article}
\usepackage[margin=2.5cm]{geometry}
\usepackage{ctex}
\usepackage{amsmath, amssymb}
\usepackage{listings}
\usepackage{xcolor}
\usepackage{tabularx}

\lstset{
    language=Python,
    basicstyle=\ttfamily\small,
    keywordstyle=\color{blue},
    commentstyle=\color{green!50!black},
    stringstyle=\color{red},
    showstringspaces=false,
    breaklines=true,
    numbers=left,
    numberstyle=\tiny\color{gray},
    frame=single
}

\title{知识点：循环神经网络 (RNN)}
\author{}
\date{}

\begin{document}
\maketitle

\section{核心定义}
\textbf{中文定义}：循环神经网络（Recurrent Neural Network, RNN）是一类用于处理序列数据（如文本、时间序列、音频）的神经网络。与传统前馈网络不同，RNN 在层间引入了反馈循环，使得模型能够保留并利用历史信息。

\textbf{英文定义}：Recurrent Neural Networks (RNNs) are a class of neural networks designed for processing sequential data. Unlike feedforward networks, RNNs feature feedback loops, allowing them to maintain an internal state (memory) to process sequences of volatile length.

\textbf{核心价值}：
RNN 的核心在于\textbf{参数共享}。在序列的每一个时间步，模型都使用相同的权重矩阵来处理输入，这使得模型能够处理任意长度的序列。

\section{核心原理与常用变体}

\subsection{Vanilla RNN (基础循环网络)}
在每个时间步 $t$，隐藏状态 $h_t$ 的更新公式为：
$$ h_t = \sigma(W_{hh} h_{t-1} + W_{xh} x_t + b_h) $$
其中 $x_t$ 是当前输入，$h_{t-1}$ 是上一时刻的记忆。
\textbf{致命缺陷}：难以处理长距离依赖。在反向传播时，长期梯度会涉及多次权重矩阵相乘，导致\textbf{梯度消失}或\textbf{梯度爆炸}。

\subsection{LSTM (Long Short-Term Memory)}
通过引入\textbf{门控机制}（Gate Mechanism）解决长期依赖问题。
\begin{enumerate}
    \item \textbf{遗忘门} ($f_t$)：决定丢弃多少旧记忆。
    \item \textbf{输入门} ($i_t$)：决定存入多少新信息。
    \item \textbf{输出门} ($o_t$)：决定输出多少当前状态。
    \item \textbf{细胞状态} ($C_t$)：贯穿所有时间步的“主干线”，梯度可在此平滑流动。
\end{enumerate}

\subsection{GRU (Gated Recurrent Unit)}
LSTM 的简化版本，将遗忘门和输入门合并为\textbf{更新门} ($z_t$)，并将细胞状态与隐藏状态合并。计算开销更小，在中小规模数据集上表现常优于 LSTM。

\section{典型架构示意图}
RNN 的展开（Unrolling）结构如下：
\begin{verbatim}
      [ h_t-1 ]      [ h_t ]      [ h_t+1 ]
          ^              ^            ^
          |              |            |
(W_hh) -> [ Cell ] ----> [ Cell ] --> [ Cell ] -> (W_hh)
          ^              ^            ^
          |              |            |
      [ x_t-1 ]      [ x_t ]      [ x_t+1 ]
\end{verbatim}

\section{数学推导：BPTT 算法}
RNN 的训练算法称为\textbf{随时间反向传播（BPTT）}。总损失 $L$ 是各时刻损失之和。对于权重 $W_{hh}$ 的梯度：
$$ \frac{\partial L}{\partial W_{hh}} = \sum_{t=1}^T \frac{\partial L_t}{\partial W_{hh}} $$
在计算 $\frac{\partial L_t}{\partial W_{hh}}$ 时，需要追溯到 $t, t-1, \dots, 1$ 时刻的所有状态，梯度项包含 $\prod_{k=1}^t \frac{\partial h_k}{\partial h_{k-1}}$。若权重矩阵特征值不等于 1，连乘会导致数值剧变。

\section{关键代码片段 (PyTorch实现)}
\begin{lstlisting}
import torch
import torch.nn as nn

class RNNModel(nn.Module):
    def __init__(self, input_size, hidden_size, output_size):
        super().__init__()
        # input_size: 词向量维度, hidden_size: 记忆单元维度
        self.rnn = nn.RNN(input_size, hidden_size, batch_first=True)
        # 也可以替换为 nn.LSTM 或 nn.GRU
        # self.rnn = nn.LSTM(input_size, hidden_size, batch_first=True)
        
        self.fc = nn.Linear(hidden_size, output_size)

    def forward(self, x):
        # x shape: (batch, seq_len, input_size)
        # out: 包含所有时间步的隐藏状态 (batch, seq_len, hidden_size)
        # hn: 最后一个时间步的隐藏状态 (1, batch, hidden_size)
        out, hn = self.rnn(x)
        
        # 通常取最后一个时间步进行分类
        last_out = out[:, -1, :]
        return self.fc(last_out)
\end{lstlisting}

\section{深度面试题与解答}
\subsection{基础概念}
\textbf{问：RNN 为什么比全连接网络更适合处理文本？}\\
答：文本具有时序相关性（前后的词有逻辑联系）。RNN 的循环结构允许信息在时间步之间传递，实现了“记忆”功能。此外，文本长度是不固定的，RNN 的参数共享机制允许它处理任意长度的序列输入。

\subsection{进阶原理}
\textbf{问：LSTM 为什么能缓解梯度消失？}\\
答：因为 LSTM 引入了“细胞状态”路径，$C_t = f_t \odot C_{t-1} + i_t \odot \tilde{C}_t$。在反向传播时，梯度在 $C_t$ 通道上主要是相加关系，而不是连乘关系。即便某些门关闭，梯度依然可以通过遗忘门路径较好地传回较远的时间步。

\subsection{工程实战}
\textbf{问：双向 RNN (Bi-RNN) 的原理是什么？}\\
答：同时运行两个 RNN，一个从前向后（正向），一个从后向前（反向）。在每个时刻 $t$，将两者的隐藏状态拼接。这使得模型在预测当前词时，既能看到上文信息，也能看到下文信息。

\subsection{坑点分析}
\textbf{问：RNN 训练时最常见的数值问题是什么？如何解决？}\\
答：\textbf{梯度爆炸}。表现为 Loss 突然变成 NaN。标准解决方法是使用\textbf{梯度裁剪（Gradient Clipping）}，即如果梯度的范数超过某个阈值，则对其进行等比例缩放，防止更新步长过大。

\section{RNN vs Transformer}

\begin{table}[h]
\centering
\begin{tabularx}{\textwidth}{|l|X|X|}
\hline
\textbf{特性} & \textbf{RNN (LSTM/GRU)} & \textbf{Transformer} \\
\hline
计算模式 & 串行（逐个单词处理） & 并行（一次性处理全句） \\
\hline
长程依赖 & 依赖记忆门（有上限） & 全局自注意力（无损） \\
\hline
硬件效率 & 较低（GPU 利用率受限） & 极高 \\
\hline
位置信息 & 隐式（结构决定顺序） & 显式（需加入位置编码） \\
\hline
\end{tabularx}
\end{table}

\section{参考文献与拓展}
\begin{enumerate}
    \item \textbf{LSTM 奠基论文}: Hochreiter, S., & Schmidhuber, J. (1997). \textit{Long short-term memory}. Neural Computation.
    \item \textbf{GRU 论文}: Cho, K., et al. (2014). \textit{Learning Phrase Representations using RNN Encoder-Decoder}.
    \item \textbf{深入浅出}: \textit{Understanding LSTMs} (Colah's blog, 最经典的博文)。
\end{enumerate}

\end{document}
